\SourceLeaf{060}
quapropter $\Omega$ erit productum ex $\varphi(u,x)$ in factores
\[
\begin{gathered}
1+\beta w\, .\,\frac{\mathrm dQ}{\mathrm dx}-\gamma w\, .\,\frac{\mathrm dQ}{\mathrm du}\\
1+\beta'w\, .\,\frac{\mathrm dQ'}{\mathrm dx}-\gamma'w\, .\,\frac{\mathrm dQ'}{\mathrm du}\\
1+\beta''w\, .\,\frac{\mathrm dQ''}{\mathrm dx}-\gamma''w\, .\,\frac{\mathrm dQ''}{\mathrm du}
\end{gathered}
\]
etc. i. e. ex $\varphi(u,x)$ in functionem integram indeterminatarum $u,x,w,\alpha,\beta,\gamma$,
$\alpha',\beta',\gamma',\alpha'',\beta'',\gamma''$ etc. Q. E. D.

\SecondArticle{17}
Theorema art. praec. manifesto applicabile est ad functionem $\zeta$, quam abhinc per
\[
f(u,x,\lambda',\lambda'',\lambda'''\text{ etc.})
\]
exhiberi supponemus, ita ut
\[
f(u+w\, .\,\tfrac{\mathrm d\zeta}{\mathrm dx},\ x-w\, .\,\tfrac{\mathrm d\zeta}{\mathrm du},\ \lambda',\lambda'',\lambda'''\text{ etc.})
\]
indefinite divisibilis evadat per $\zeta$: quotientem, qui erit functio integra indeterminatarum $u,x,w,a,b,c$ etc., symmetrica respectu ipsarum $a,b,c$ etc., exhibebimus per
\[
\psi(u,x,w,\lambda',\lambda'',\lambda'''\text{ etc.})
\]
Hinc concludimus, fieri etiam identice
\[
\begin{aligned}
&f(u+w\, .\,\tfrac{\mathrm dz}{\mathrm dx},\ x-w\, .\,\tfrac{\mathrm dz}{\mathrm du},\ l',l'',l'''\text{ etc.})\\
&=z\psi(u,x,w,l',l'',l'''\text{ etc.})
\end{aligned}
\]
nec non
\[
\begin{aligned}
&f(u+w\, .\,\tfrac{\mathrm dZ}{\mathrm dx},\ x-w\, .\,\tfrac{\mathrm dZ}{\mathrm du},\ L',L'',L'''\text{ etc.})\\
&=Z\psi(u,x,w,L',L'',L'''\text{ etc.})
\end{aligned}
\]
Quodsi itaque functionem $Z$ simpliciter exhibemus per $F(u,x)$, ita ut habeatur
\[
f(u,x,L',L'',L'''\text{ etc.})=F(u,x)
\]
erit identice
\[
\begin{aligned}
&F(u+w\, .\,\tfrac{\mathrm dZ}{\mathrm dx},\ x-w\, .\,\tfrac{\mathrm dZ}{\mathrm du})\\
&=Z\psi(u,x,w,L',L'',L'''\text{ etc.})
\end{aligned}
\]

\SourceLeaf{061}
\SecondArticle{18}
Si itaque e valoribus determinatis ipsarum $u,x$, puta ex $u=U,x=X$,
prodire supponimus
\[
\frac{\mathrm dZ}{\mathrm dx}=X',\quad\frac{\mathrm dZ}{\mathrm du}=U'
\]
erit identice
\[
\begin{aligned}
&F(U+wX',\ X-wU')\\
&=F(U,X)\, .\,\psi(U,X,w,L',L'',L'''\text{ etc.})
\end{aligned}
\]
Quoties $U'$ non evanescit, statuere licebit
\[
w=\frac{X-x}{U'}
\]
unde emergit
\[
\begin{aligned}
&F(U+\tfrac{XX'}{U'}-\tfrac{X'x}{U'},\ x)\\
&=F(U,X)\, .\,\psi(U,X,\tfrac{X-x}{U'},L',L'',L'''\text{ etc.})
\end{aligned}
\]
quod etiam ita enunciare licet:

Si in functione $Z$ statuitur $u=U+\frac{XX'}{U'}-\frac{X'x}{U'}$, transit ea in
\[
F(U,X)\, .\,\psi(U,X,\tfrac{X-x}{U'},L',L'',L'''\text{ etc.})
\]

\SecondArticle{19}
Quum in casu eo, ubi non est $P=0$, determinans functionis $Z$ sit
functio indeterminatae $x$ per se non evanescens, manifesto multitudo valorum determinatorum ipsius $x$, per quos hic determinans valorem 0 nancisci potest, erit
numerus finitus, ita ut infinite multi valores determinati ipsius $x$ assignari possint, qui determinanti illi valorem a 0 diversum concilient. Sit $X$ talis valor
ipsius $x$ (quem insuper \textit{realem} supponere licet). Erit itaque determinans functionis $F(u,X)$ non $=0$, unde sequitur, per theorema II. art. 6, functiones
\[
F(u,X)\quad\text{et}\quad\frac{\mathrm dF(u,X)}{\mathrm du}
\]
habere non posse divisorem ullum communem. Supponamus porro, exstare aliquem
valorem determinatum ipsius $u$, puta $U$ (sive realis sit, sive imaginarius i. e. sub
forma $g+h\BareRoot{-1}$ contentus), qui reddat $F(u,X)=0$, i. e. esse $F(U,X)=0$.
Erit itaque $u-U$ factor indefinitus functionis $F(u,X)$, et proin functio
$\frac{\mathrm dF(u,X)}{\mathrm du}$ certo per $u-U$ non divisibilis. Supponendo itaque, hanc functionem
